\documentclass[reqno]{amsart}
\usepackage[T1]{fontenc}
\usepackage{amsmath,amssymb,amsthm}
\usepackage{xr-hyper}
\usepackage[hidelinks]{hyperref}
\externaldocument[st-]{20261124T214500--supersingular-theta-divisibility__workbook}
\externaldocument[ip-]{20261124T214600--the-isotropic-plane__workbook}
\externaldocument[ba-]{20261126T030200--the-bound-is-attained__workbook}
\externaldocument[da-]{20261126T030300--the-data__workbook}
\externaldocument[ob-]{20261126T040100--spaces-forms-and-twists__workbook}
\externaldocument[dl-]{20261126T030500--divisibility-of-the-composite-lift__workbook}
\externaldocument[d2-]{20261126T050200--one-lift-at-a-small-non-split-prime__workbook}
\externaldocument[c3-]{20261126T050600--the-exchange-identity__workbook}
\externaldocument[col-]{20261126T030800--the-cost-of-one-lift__workbook}
\externaldocument[l3-]{20261126T050900--hecke-values-as-norms-of-lifts__workbook}
\externaldocument[l5-]{20261126T051100--test-forms-and-the-height-of-a-hecke-ratio__workbook}

\newcounter{chap}
% Chapter numbers, generated from 20261125T012000--grh-workbook__guide.org
% by scripts/make-chapters.py.  Do not edit: change the order of the
% index and rerun the script (./preview-site.sh and CI do this).
%
% Usage in a sheet, after \newcounter{chap}:
%   % Chapter numbers, generated from 20261125T012000--grh-workbook__guide.org
% by scripts/make-chapters.py.  Do not edit: change the order of the
% index and rerun the script (./preview-site.sh and CI do this).
%
% Usage in a sheet, after \newcounter{chap}:
%   % Chapter numbers, generated from 20261125T012000--grh-workbook__guide.org
% by scripts/make-chapters.py.  Do not edit: change the order of the
% index and rerun the script (./preview-site.sh and CI do this).
%
% Usage in a sheet, after \newcounter{chap}:
%   \input{chapters}\chapnumber{<sheet ID>}
\makeatletter
\providecommand\chapters@set[2]{\@namedef{chapters@#1}{#2}}
\chapters@set{20261125T001000}{1}% zeros-to-ramp-sums
\chapters@set{20261125T001100}{2}% inert-primes-and-estimate-b
\chapters@set{20261125T005000}{3}% galois-fourier-and-averaging
\chapters@set{20261125T005500}{4}% columns-from-a-sphere
\chapters@set{20261125T000000}{5}% slope-method
\chapters@set{20261126T030100}{6}% a-twist-for-each-prime
\chapters@set{20261125T003000}{7}% the-chain-and-the-square-root
\chapters@set{20261125T010000}{8}% rank-lemma-and-zero-estimate
\chapters@set{20261125T010500}{9}% poincare-series-and-near-orthogonality
\chapters@set{20261124T223000}{10}% elliptic-curves-for-the-proof
\chapters@set{20261125T009500}{11}% fock-model-and-theta-intertwiners
\chapters@set{20261124T223100}{12}% complex-multiplication-and-inert-reduction
\chapters@set{20261125T008000}{13}% one-cm-determinant
\chapters@set{20261124T214500}{14}% supersingular-theta-divisibility
\chapters@set{20261124T214600}{15}% the-isotropic-plane
\chapters@set{20261125T013700}{16}% pairings-on-a-formal-group-of-height-two
\chapters@set{20261125T013500}{17}% hecke-nondegeneracy-on-the-tree
\chapters@set{20261126T030200}{18}% the-bound-is-attained
\chapters@set{20261125T009000}{19}% kronecker-limit-and-chowla-selberg
\chapters@set{20261125T002000}{20}% definite-theta-lifts
\chapters@set{20261125T015500}{21}% siegel-weil-in-the-smallest-case
\chapters@set{20261125T002100}{22}% definite-theta-lifts-and-norms
\chapters@set{20261125T016500}{23}% hecke-operators-and-theta-lifts
\chapters@set{20261125T016700}{24}% divisibility-along-the-chain
\chapters@set{20261125T004000}{25}% denominators-and-sha
\chapters@set{20261125T004100}{26}% tate-twists-and-sha
\chapters@set{20261126T020430}{27}% the-minus-part-by-gauss-sums
\chapters@set{20261125T014100}{28}% congruence-modules
\chapters@set{20261125T014500}{29}% eisenstein-congruences-and-selmer-classes
\chapters@set{20261125T014700}{30}% pseudocharacters-and-the-length-bound
\chapters@set{20261125T011000}{31}% galois-representations-black-box
\chapters@set{20261126T030300}{32}% the-data
\chapters@set{20261126T020210}{33}% bessel-poincare-series-and-their-gram-operators
\chapters@set{20261126T030400}{34}% the-seesaw-identity
\chapters@set{20261126T030500}{35}% divisibility-of-the-composite-lift
\chapters@set{20261126T030600}{36}% theorem-a-short-forms
\chapters@set{20261126T030700}{37}% the-determinant-of-a-chain
\chapters@set{20261126T030800}{38}% the-cost-of-one-lift
\chapters@set{20261126T030900}{39}% denominators-of-projections-to-lifts
\chapters@set{20261126T031000}{40}% twists-with-few-congruences
\chapters@set{20261126T031100}{41}% theorem-c-and-the-main-inequality
\providecommand\chapnumber[1]{%
  \@ifundefined{chapters@#1}%
    {\PackageError{chapters}{Sheet #1 is not listed in the index}%
      {List it in 20261125T012000--grh-workbook__guide.org\MessageBreak
       and run scripts/make-chapters.py.}}%
    {\edef\thechap{\@nameuse{chapters@#1}}}}
\makeatother
\chapnumber{<sheet ID>}
\makeatletter
\providecommand\chapters@set[2]{\@namedef{chapters@#1}{#2}}
\chapters@set{20261125T001000}{1}% zeros-to-ramp-sums
\chapters@set{20261125T001100}{2}% inert-primes-and-estimate-b
\chapters@set{20261125T005000}{3}% galois-fourier-and-averaging
\chapters@set{20261125T005500}{4}% columns-from-a-sphere
\chapters@set{20261125T000000}{5}% slope-method
\chapters@set{20261126T030100}{6}% a-twist-for-each-prime
\chapters@set{20261125T003000}{7}% the-chain-and-the-square-root
\chapters@set{20261125T010000}{8}% rank-lemma-and-zero-estimate
\chapters@set{20261125T010500}{9}% poincare-series-and-near-orthogonality
\chapters@set{20261124T223000}{10}% elliptic-curves-for-the-proof
\chapters@set{20261125T009500}{11}% fock-model-and-theta-intertwiners
\chapters@set{20261124T223100}{12}% complex-multiplication-and-inert-reduction
\chapters@set{20261125T008000}{13}% one-cm-determinant
\chapters@set{20261124T214500}{14}% supersingular-theta-divisibility
\chapters@set{20261124T214600}{15}% the-isotropic-plane
\chapters@set{20261125T013700}{16}% pairings-on-a-formal-group-of-height-two
\chapters@set{20261125T013500}{17}% hecke-nondegeneracy-on-the-tree
\chapters@set{20261126T030200}{18}% the-bound-is-attained
\chapters@set{20261125T009000}{19}% kronecker-limit-and-chowla-selberg
\chapters@set{20261125T002000}{20}% definite-theta-lifts
\chapters@set{20261125T015500}{21}% siegel-weil-in-the-smallest-case
\chapters@set{20261125T002100}{22}% definite-theta-lifts-and-norms
\chapters@set{20261125T016500}{23}% hecke-operators-and-theta-lifts
\chapters@set{20261125T016700}{24}% divisibility-along-the-chain
\chapters@set{20261125T004000}{25}% denominators-and-sha
\chapters@set{20261125T004100}{26}% tate-twists-and-sha
\chapters@set{20261126T020430}{27}% the-minus-part-by-gauss-sums
\chapters@set{20261125T014100}{28}% congruence-modules
\chapters@set{20261125T014500}{29}% eisenstein-congruences-and-selmer-classes
\chapters@set{20261125T014700}{30}% pseudocharacters-and-the-length-bound
\chapters@set{20261125T011000}{31}% galois-representations-black-box
\chapters@set{20261126T030300}{32}% the-data
\chapters@set{20261126T020210}{33}% bessel-poincare-series-and-their-gram-operators
\chapters@set{20261126T030400}{34}% the-seesaw-identity
\chapters@set{20261126T030500}{35}% divisibility-of-the-composite-lift
\chapters@set{20261126T030600}{36}% theorem-a-short-forms
\chapters@set{20261126T030700}{37}% the-determinant-of-a-chain
\chapters@set{20261126T030800}{38}% the-cost-of-one-lift
\chapters@set{20261126T030900}{39}% denominators-of-projections-to-lifts
\chapters@set{20261126T031000}{40}% twists-with-few-congruences
\chapters@set{20261126T031100}{41}% theorem-c-and-the-main-inequality
\providecommand\chapnumber[1]{%
  \@ifundefined{chapters@#1}%
    {\PackageError{chapters}{Sheet #1 is not listed in the index}%
      {List it in 20261125T012000--grh-workbook__guide.org\MessageBreak
       and run scripts/make-chapters.py.}}%
    {\edef\thechap{\@nameuse{chapters@#1}}}}
\makeatother
\chapnumber{<sheet ID>}
\makeatletter
\providecommand\chapters@set[2]{\@namedef{chapters@#1}{#2}}
\chapters@set{20261125T001000}{1}% zeros-to-ramp-sums
\chapters@set{20261125T001100}{2}% inert-primes-and-estimate-b
\chapters@set{20261125T005000}{3}% galois-fourier-and-averaging
\chapters@set{20261125T005500}{4}% columns-from-a-sphere
\chapters@set{20261125T000000}{5}% slope-method
\chapters@set{20261126T030100}{6}% a-twist-for-each-prime
\chapters@set{20261125T003000}{7}% the-chain-and-the-square-root
\chapters@set{20261125T010000}{8}% rank-lemma-and-zero-estimate
\chapters@set{20261125T010500}{9}% poincare-series-and-near-orthogonality
\chapters@set{20261124T223000}{10}% elliptic-curves-for-the-proof
\chapters@set{20261125T009500}{11}% fock-model-and-theta-intertwiners
\chapters@set{20261124T223100}{12}% complex-multiplication-and-inert-reduction
\chapters@set{20261125T008000}{13}% one-cm-determinant
\chapters@set{20261124T214500}{14}% supersingular-theta-divisibility
\chapters@set{20261124T214600}{15}% the-isotropic-plane
\chapters@set{20261125T013700}{16}% pairings-on-a-formal-group-of-height-two
\chapters@set{20261125T013500}{17}% hecke-nondegeneracy-on-the-tree
\chapters@set{20261126T030200}{18}% the-bound-is-attained
\chapters@set{20261125T009000}{19}% kronecker-limit-and-chowla-selberg
\chapters@set{20261125T002000}{20}% definite-theta-lifts
\chapters@set{20261125T015500}{21}% siegel-weil-in-the-smallest-case
\chapters@set{20261125T002100}{22}% definite-theta-lifts-and-norms
\chapters@set{20261125T016500}{23}% hecke-operators-and-theta-lifts
\chapters@set{20261125T016700}{24}% divisibility-along-the-chain
\chapters@set{20261125T004000}{25}% denominators-and-sha
\chapters@set{20261125T004100}{26}% tate-twists-and-sha
\chapters@set{20261126T020430}{27}% the-minus-part-by-gauss-sums
\chapters@set{20261125T014100}{28}% congruence-modules
\chapters@set{20261125T014500}{29}% eisenstein-congruences-and-selmer-classes
\chapters@set{20261125T014700}{30}% pseudocharacters-and-the-length-bound
\chapters@set{20261125T011000}{31}% galois-representations-black-box
\chapters@set{20261126T030300}{32}% the-data
\chapters@set{20261126T020210}{33}% bessel-poincare-series-and-their-gram-operators
\chapters@set{20261126T030400}{34}% the-seesaw-identity
\chapters@set{20261126T030500}{35}% divisibility-of-the-composite-lift
\chapters@set{20261126T030600}{36}% theorem-a-short-forms
\chapters@set{20261126T030700}{37}% the-determinant-of-a-chain
\chapters@set{20261126T030800}{38}% the-cost-of-one-lift
\chapters@set{20261126T030900}{39}% denominators-of-projections-to-lifts
\chapters@set{20261126T031000}{40}% twists-with-few-congruences
\chapters@set{20261126T031100}{41}% theorem-c-and-the-main-inequality
\providecommand\chapnumber[1]{%
  \@ifundefined{chapters@#1}%
    {\PackageError{chapters}{Sheet #1 is not listed in the index}%
      {List it in 20261125T012000--grh-workbook__guide.org\MessageBreak
       and run scripts/make-chapters.py.}}%
    {\edef\thechap{\@nameuse{chapters@#1}}}}
\makeatother

\ifcsname chapters@20261126T051000\endcsname
  \chapnumber{20261126T051000}
\else
  \renewcommand{\thechap}{?}
\fi
\renewcommand{\thesection}{\thechap.\arabic{section}}
\theoremstyle{definition}
\newtheorem{problem}{Problem}[section]
\newtheorem{definition}{Definition}[section]
\theoremstyle{plain}
\newtheorem{theorem}{Theorem}[section]
\newcommand{\Z}{\mathbb{Z}}
\newcommand{\Q}{\mathbb{Q}}
\newcommand{\C}{\mathbb{C}}
\newcommand{\N}{\mathbb{N}}
\newcommand{\ii}{\mathrm{i}}
\newcommand{\Oc}{\mathcal{O}}
\newcommand{\sll}{\operatorname{sl}}
\newcommand{\val}{\operatorname{val}}
\newcommand{\lv}{\operatorname{lv}}
\newcommand{\Aut}{\operatorname{Aut}}
\newcommand{\Hom}{\operatorname{Hom}}
\newcommand{\sfp}{s_{\mathrm{fin},p}}
\newcommand{\stot}{s^{\mathrm{tot}}_{\mathrm{fin},p}}
\newcommand{\cvac}{c_{\mathrm{vac}}}
\newcommand{\catt}{c_{\mathrm{att}}}
\newcommand{\cjet}{c_{\mathrm{jet}}}
\newcommand{\cper}{c_{\mathrm{per}}}
\newcommand{\Jet}{\mathrm{Jet}}
\newcommand{\Th}{\mathrm{Th}}
\newcommand{\tth}{\mathrm{th}}
\newcommand{\thfr}{\mathrm{th}_{\mathrm{fr}}}
\newcommand{\hb}{\mathbf h}
\newcommand{\oo}{\mathfrak o}
\setlength{\emergencystretch}{3em}

\title{The bound is attained at many places}
\author{}
\date{November 26, 2026}

\begin{document}
\maketitle
\refstepcounter{chap}\label{chap}

\paragraph{Prerequisites.}
Chapter~\ref{ba-chap}: the least valuation \(\mu(n)\) on the lemniscate, Theorems~\ref{ba-thm:rate} and~\ref{ba-thm:katz} with the series \(y(z)\) and the ring \(\oo[[y]]\) of Section~\ref{ba-sec:general}, and Problems~\ref{ba-pr:monomial}, \ref{ba-pr:gamma-cases}, \ref{ba-pr:many-variables}, \ref{ba-pr:prime-sum-rate}, \ref{ba-pr:product} and~\ref{ba-pr:near-pairs}.  The moments \(t(a,w)\) of Section~\ref{st-sec:moments} and Theorem~\ref{ip-thm:main}.  The data (Definition~\ref{da-def:data}, the constant \(v_0\) of Section~\ref{da-sec:data}), the integral vacuum (Section~\ref{ob-sec:vacuum}), the local prime sums \(\sfp\), the numbers \(c_p\) and the sets \(I_p\) (Section~\ref{dl-sec:gain}), and the jet series of a datum with its scale \(\cjet\) (Definition~\ref{d2-def:jet-series}).  Lagrange's four-square theorem; the ultrametric inequality in \(\C_p\), whose valuation \(\val\) has \(\val(p) = 1\).

\medskip
Chapter~\ref{col-chap} bounds the cost of one lift using the test forms of Chapter~\ref{l5-chap}.  At a prime \(p\) prime to the norm of the twist, a test form uses a Taylor coefficient on twelve copies of the variables; multiplied by \(\cvac^{12}\prod_\tau\lambda(\tau)!\) and divided by one \(\catt\) independent of \(p\), it must have valuation at most the divisibility \(\stot\) predicted by the prime sum.  How can one normalization serve all primes, when the partners \([-p^e]\tau\) of a place lower the prediction?

\section{One curve, one partner}\label{sec:lemniscate}

Recall \(\mu(n) \le \gamma_p n\), \(\gamma_p = p/(p^2-1)\), for \(p \equiv 3 \pmod 4\) (Theorem~\ref{ba-thm:rate}).

\begin{problem}\label{pr:lemniscate-least}
For \(p = 3\): \(\mu(3), \mu(9), \mu(11), \mu(81) = 1, 3, 4, 30\), each below \(3n/8\).
\end{problem}

For \(p \equiv 3 \pmod 4\), \(w \ge 0\) and \(G \in \Z_p[\ii][[X_1, X_2]]\) with \(X_1 = \sll s\), \(X_2 = \sll t\), the valuation of \(\partial_s^{pw}\partial_t^w G(0,0)\) is at least \(\mu(pw) + \mu(w)\) (Problem~\ref{ba-pr:product}).  A non-integral series does better.

\begin{problem}\label{pr:coupling-exp}
\(\partial_s^{pw}\partial_t^{w} e^{s^pt/p}\big|_{(0,0)} = (pw)!/p^w\), of valuation \(v_p(w!)\).  For \(p = 3\), \(w = 2\) this is \(80\), a unit, while for every \(G \in \Z_3[\ii][[X_1, X_2]]\), with \(X_1 = \sll s\), \(X_2 = \sll t\), the valuation of \(\partial_s^6\partial_t^2 G(0,0)\) is at least \(2\).
\end{problem}

For a prime \(p \equiv 3 \pmod 8\) and integers \(a, w \ge 0\), Theorem~\ref{ip-thm:main} gives in particular \(v_p(t(a,w)) \ge \sum_{i \ge 1,\ i \text{ odd}} (\lfloor a/p^i\rfloor - w)_+\) for the moments of Section~\ref{st-sec:moments}, where \(w\) is the order of the derivative in the partner direction.

\begin{problem}\label{pr:one-partner}
\(v_{11}(t(27,1)) = v_{11}(t(38,2)) = v_{11}(t(49,3)) = 1\), and in each case the bound of Theorem~\ref{ip-thm:main} at \(p = 11\) is \(1\): in these three cases the theta function itself attains the bound.
\end{problem}

\begin{problem}\label{pr:not-attained}
At \(p = 3\): \(v_3(t(a,1)) = 6, 6, 9\) for \(a = 11, 19, 27\), while the bound of Theorem~\ref{ip-thm:main} is \(2, 5, 8\).
\end{problem}

So a fixed theta function can miss.  For each \(p\) and \(\eta\) the proof chooses twelve sections, one on each copy of the variables \((z_\tau)_{\tau\in\Phi}\).

\section{What the partners take away}\label{sec:partners}

From now on we use the data (Definition~\ref{da-def:data}).  Fix a prime \(p\) and a ring homomorphism \(\eta\colon\C\to\C_p\).  Then \(c_p\), \(I_p\) and the local prime sum \(\sfp(y,\sigma) = c_p \sum_{e \in I_p} (y/p^e - \kappa([-p^e]\sigma))_+\) are those of Section~\ref{dl-sec:gain}.  Let \(\gamma_p := c_p\sum_{e \in I_p} p^{-e}\) (values in Problem~\ref{ba-pr:gamma-cases}).  The \emph{total local prime sum} is
\begin{equation}\label{eq:stot}
  \stot := \sum_{\tau \in \Phi} \sfp(\kappa(\tau), \tau) .
\end{equation}
By Problem~\ref{ba-pr:prime-sum-rate}, \(\stot \le \gamma_p\sum_\tau \kappa(\tau)\).  The difference is what the partners take away.

\begin{problem}\label{pr:partners}
\(\gamma_p \sum_{\tau} \kappa(\tau) - \stot = c_p \sum_{\tau \in \Phi} \sum_{e \in I_p} \min\bigl(\kappa(\tau)/p^e,\ \kappa([-p^e]\tau)\bigr)\).
\end{problem}

At an inert prime the partners enter the theta function through couplings.  The proof uses a fixed curve: the lattice \(\Lambda = \Omega_0\Oc_K\), with \(\Omega_0\) as in Definition~\ref{d2-def:jet-series}, with \(t\), \(a_1\) and the series \(y\) as in Section~\ref{ba-sec:general}; let \(\operatorname{area}(\Lambda)\) be the area of \(\C/\Lambda\).  Let \(\sigma_\Lambda\) be the odd entire function with \(\sigma_\Lambda'(0) = 1\) and \(\sigma_\Lambda(z + l) = \epsilon_l \exp\bigl(\pi(\bar l z + \lvert l\rvert^2/2)/\operatorname{area}(\Lambda)\bigr)\sigma_\Lambda(z)\) for \(l \in \Lambda\), where \(\epsilon_l = 1\) on \(2\Lambda\) and \(-1\) elsewhere.  In this sheet let
\[
  \tth(z) := e^{-\pi z\bar t/\operatorname{area}(\Lambda)}\, \frac{\sigma_\Lambda(z + t)}{\sigma_\Lambda(t)} ,
\]
We read its Taylor series in \(\C_p\) through \(\eta\).  For linear forms \(\mu_\beta\) on \(\C_p^{\Phi}\) and integers \(c_\beta \ge 0\) (\(\beta\) in a finite set) put \(\thfr(z) := \prod_\beta \tth(\mu_\beta z)^{c_\beta}\), and for odd \(e\) let \(\Gamma_e := \sum_{\tau} z_\tau^{p^e} z_{[-p^e]\tau}\) be the coupling of each place with its partner.

\paragraph{Normal form.}  Let \(p\) be inert in \(K\), and let the forms be those of a tight frame (Definition~\ref{def:twelve} below).  The proof shows that there are \(\varpi \in \C_p\) with \(\val\varpi = 1/(p^2-1)\), integral \(\mathfrak l_{e,j}\) (\(e\) odd, \(j \ge 0\)) with \(\mathfrak l_{e,0}\) a unit, and an integral series \(W\) in the \(z_\tau\), such that
\[
  \thfr(\varpi z) = W(z) \exp\Bigl(\frac{\varpi^{p+1}}{p} \sum_{e\ \mathrm{odd}} \sum_{j \ge 0} \mathfrak l_{e,j}\, \Gamma_e^{\,p^j}\Bigr) .
\]
\emph{Status}: read from the proof.

\begin{problem}\label{pr:coupled-pair}
A two-variable model: \(K = \Q(\sqrt{-7})\), \(p = 13\), variables \(z_1, z_2\), \(\val\varpi = 1/168\) and \(\mathfrak l_{1,0}\) a unit.  Then \(13\) is inert in \(K\), \(\val(\varpi^{14}/13) = -11/12\), and the coefficient of \(z_1^{13}z_2\) in \((\varpi^{14}/13)\mathfrak l_{1,0}(z_1/\varpi)^{13}(z_2/\varpi)\) is \(\mathfrak l_{1,0}/13\), of valuation \(-1\); times \(13!\,1!\) it has valuation \(0\), whereas the rate of one variable in degree \(14\) is \(\gamma_{13}\cdot 14 = 13/12\).
\end{problem}

\section{Six copies along an isotropic vector}\label{sec:six}

\begin{definition}\label{def:twelve}
A \emph{tight frame} of length \(M\) consists of elements \(w_\beta \in E\) and integers \(c_\beta \ge 0\), \(\beta\) in a finite set, with \(\sum_\beta c_\beta\,\tau'(w_\beta)\overline{\tau(w_\beta)} = M\delta_{\tau'\tau}\) for \(\tau, \tau' \in \Phi\), each \(w_\beta\) lying in the lattice \(\mathfrak a_0 := E \cap \alpha_0\hat\Oc_E\), \(\alpha_0\) the finite idele of the \emph{Weil data} fixed with \(E\) in the proof.  At \(\eta\) its forms are \(\mu_\beta(z) := \sum_{\tau\in\Phi}\eta\bigl(\cper\overline{\tau(w_\beta)}\bigr)z_\tau\), where \(\cper\) is the period factor of Definition~\ref{d2-def:jet-series} (not the \(\Omega\) of Chapter~\ref{da-chap}).  The \emph{twelve copies} are the variables \(z_{\tau,i}\), \(\tau \in \Phi\), \(1 \le i \le 12\), one coordinate of \(\Hom(E, E^{12})\) at each infinite place and each of the twelve basis vectors; a multi-index on them is written \(\hb = (h_{\tau,i})\), and \(n\) copies are the variables \(z_{\tau,i}\) with \(i \le n\).
\end{definition}

\paragraph{Tight frames.}  There is a tight frame of length \(M\) with \(c_\beta \ge 6fb\) for every \(\beta\) such that for every prime \(p\) and every \(\eta\colon\C\to\C_p\) some form \(\mu_{\beta_0}\) has all its coefficients units.  \emph{Status}: read from the proof.  Fix one; write \(\oo\) for the valuation ring of \(\C_p\), and \(\oo[[y_{\mathrm{fr}}]]\) for the ring of series \(G\bigl((y(\mu_\beta z))_\beta\bigr)\), \(G\) a power series over \(\oo\) in one variable for each \(\beta\).

\begin{definition}\label{def:lv}
Let \(n \in \N\), \(\nu\) a multi-index on \(\Phi\) and \(\rho \in \mathbb R\).  We write \(\lv_n(\nu) \le \rho\) if there are \(H_1, \dots, H_n \in \oo[[y_{\mathrm{fr}}]]\) and multi-indices \(\nu_1 + \dots + \nu_n = \nu\) with
\(\sum_{i=1}^n \val\bigl([z^{\nu_i}](H_i\,\thfr)\bigr) \le \rho\).
\end{definition}

The next problems use six scalings to make the product of the theta factors a unit of \(\oo[[y_{\mathrm{fr}}]]\).

\paragraph{Isotropic scalings.}  Write \(\tth = \exp(-g)\) with \(g(0) = 0\).  For every \(\delta \in \Z_p\) the series \(\exp\bigl(\delta^2 g(z) - g(\delta z)\bigr)\) lies in \(\oo[[y]]\).  \emph{Status}: read from the proof.  For \(0 \ne \alpha \in \Oc_K\), multiplication by \(\alpha\) is an isogeny of degree \(\alpha\bar\alpha\) of \(\C/\Lambda\) (complex multiplication; Silverman, \emph{Advanced Topics in the Arithmetic of Elliptic Curves}, Ch.~II).  For every \(\alpha \in \Oc_K\) the proof shows \(\tth(\alpha z) = \tth(z)^{\alpha\bar\alpha}v\) with \(v \in \oo[[y]]\), that is, \(\exp(\alpha\bar\alpha g(z) - g(\alpha z)) \in \oo[[y]]\); it then passes to \(p\)-adic limits of the pairs \((\alpha, \alpha\bar\alpha)\), among which are the pairs \((\delta, \delta^2)\), \(\delta \in \Z_p\).

\begin{problem}\label{pr:isotropic}
\(1 + 70^2 = 29\cdot 13^2\); \((1, 70, 0, 0, 0, 0)\) is isotropic modulo \(13^2\); there is \(\delta_2 \in \Z_{13}\) with \(\delta_2 \equiv 70 \pmod{169}\) and \(\delta_2^2 = -1\).  For every prime \(p\) there are \(\delta_2, \dots, \delta_5 \in \Z_p\) such that \(\delta := (1, \delta_2, \delta_3, \delta_4, \delta_5, 0)\) has \(\sum_{i=1}^6 \delta_i^2 = 0\).
\end{problem}

\begin{problem}\label{pr:unit-inverse}
If \(\delta \in \Z_p^6\) and \(\sum_i\delta_i^2 = 0\), then \(\prod_{i}\tth(\delta_i z)\) lies in \(\oo[[y]]\) with constant term \(1\), and \(\prod_i \thfr(\delta_i z)\) is a unit of \(\oo[[y_{\mathrm{fr}}]]\).
\end{problem}

\begin{problem}\label{pr:six-unit}
If \(\mu_{\beta_0}\) has unit coefficients, then \(\lv_6(\nu) \le \gamma_p\lvert\nu\rvert - \sum_{\tau}\val(\nu_\tau!)\) for every \(\nu\).  For \(\nu = (3,6)\) on two places and \(p = 3\): \(\binom93 = 84\), \(v_3(9!) = 4\) and \(v_3(3!\,6!) = 3\).
\end{problem}

Two more six-copy bounds are read from the proof.  At an inert prime \(p\), integral combinations of the forms give each coordinate times a unit, and then \(\lv_6(\nu) \le 0\) for every \(\nu\).  At an inert prime \(p > 2\), from the normal form: let \(\Phi_0 \subseteq \Phi\), odd positive integers \(e_\tau\) for \(\tau \in \Phi_0\) such that no \((\tau_1, e')\) with \(\tau_1 \in \Phi_0\), \(e' \le e_\tau\) an odd positive integer and \((\tau_1,e') \ne (\tau,e_\tau)\) has \([-p^{e'}]\tau_1 = [-p^{e_\tau}]\tau\), integers \(j_\tau \ge 0\), \(\nu' := \sum_{\tau\in\Phi_0} p^{e_\tau}j_\tau[\tau]\), \(\nu'' := \sum_{\tau\in\Phi_0} j_\tau[[-p^{e_\tau}]\tau]\) and \(J := \sum_{\tau\in\Phi_0} j_\tau\); then
\[
  \lv_6(\nu' + \nu'') \le \frac{J}{p-1} - J - \frac{\lvert\nu' + \nu''\rvert}{p^2-1} - \val(\nu''!) ,
\]
where \([\tau]\) is the multi-index with a single \(1\) at \(\tau\) and \(\nu''! := \prod_\tau \nu''_\tau!\).

\section{Near pairs}\label{sec:near}

We now describe the pairs used at non-split primes \(p \le k\).  Let \(l_0(s)\) be as in Definition~\ref{da-def:data}.

\begin{definition}\label{def:near}
Let \(p\) be a prime.  A pair \((\tau, e)\) with \(\tau \in \Phi\) is \emph{near} at \(p\) if \(e \in I_p\), \(2e < l_0(s)\), \(\tau \in \Phi_b\) and \([-p^e]\tau \in \Phi_b\); its \emph{start} is \(\tau\) and its \emph{end} is \([-p^e]\tau\).
\end{definition}

A model: the group \((\Z/13)^\times\) in place of \(\Phi\), \(p = 3\) treated as inert, the cutoff \(2e < 3\), and the heavy set \(\{1, 10\}\) in place of \(\Phi_b\).

\begin{problem}\label{pr:model-edges}
The iterates of \([-3]\) of orders \(0\) to \(4\) from \(1\) and from \(10\) are \((1, 10, 9, 12, 3)\) and \((10, 9, 12, 3, 4)\); the only near pair is \((1, 1)\), with end \(10\).
\end{problem}

\begin{problem}\label{pr:model-sum}
With weight \(27\) on \(\{1, 10\}\) and \(1\) elsewhere, \(\sum_{\tau}\sum_{e\ \mathrm{odd}} (\kappa(\tau)/3^e - \kappa([-3^e]\tau))_+ = 8\), the only nonzero term being \(\tau = 10\), \(e = 1\); the near pair \((1,1)\) contributes \(0\).
\end{problem}

\begin{problem}\label{pr:near-unique}
Let \(p \le k\) be a prime with \(\chi_K(p) \ne 1\).  (a) Two distinct near pairs at \(p\) have distinct starts and distinct ends, and the end of a near pair is never the start of one.  (b) If \(p = d\), there are no near pairs.  Hence the near pairs at a prime \(p \le k\) satisfy the hypothesis on \(\Phi_0\) in Section~\ref{sec:six}.
\end{problem}

\section{Twelve copies and the theorem}\label{sec:twelve}

On twelve copies the proof combines the two kinds of six-copy bounds: at a split prime, Problem~\ref{pr:six-unit} on six copies and the constant coefficient in degree zero on the other six; at an inert \(p > k\), every place coupled with its partner for \(e = 1\) on six copies and the span bound on six; at a non-split \(p \le k\), the near pairs on six copies and Problem~\ref{pr:six-unit} on six.  For \(\lambda\colon\Phi\to\N\) let \(u := \lceil 6b/k\rceil\) and
\[
  \Pi_\tau := \prod_{i=1}^{u}\bigl((\lambda(\tau) - u)_+ + i\bigr),
\]
a product of \(u\) consecutive integers; the reserved \(u\) pays for the rounding.  The proof shows: if \(\kappa(\tau) \le \lambda(\tau) \le \kappa(\tau) + x^{42}\) for every \(\tau\), then at every prime \(p\) and every \(\eta\)
\[
  \lv_{12}(\lambda) \le \stot + 2\sum_\tau \val(\Pi_\tau) - \sum_\tau \val(\lambda(\tau)!) .
\]
\emph{Status}: read from the proof.  The loss added to \(\stot - \sum_\tau\val(\lambda(\tau)!)\) is zero at split primes, \(\sum_\tau\val(\Pi_\tau)\) at non-split \(p \le k\), and \(2\sum_\tau\val(\Pi_\tau)\) at inert \(p > k\).

\begin{problem}\label{pr:loss-product}
\(12\cdot13\cdot14\cdot15 = 32760\) has \(v_{13} = 1\).  A product of \(u\) consecutive integers has \(v_p \ge \lfloor u/p\rfloor\).
\end{problem}

\begin{problem}\label{pr:monomial-witness}
If \(\lv_n(\nu) \le \rho\), then the \(H_i\) in Definition~\ref{def:lv} may be taken to be monomials \(\prod_\beta y(\mu_\beta z)^{e_{i,\beta}}\) with \(\sum_\beta e_{i,\beta} \le \lvert\nu\rvert\).
\end{problem}

For a modulus \(\mathfrak m\) of \(E\) (Definition~\ref{c3-def:modulus}) with ideal \((\pi_M)\) and \(y_0 \in (\Oc_E/\mathfrak m)^{12}\), let \(\Th_{y_0}\) be the theta function (the \(\mathrm{th}_\varphi\) of Definition~\ref{d2-def:jet-series}; in this sheet \(\tth\) is the theta function of the curve) of the datum on \(\Hom(E, E^{12})\) with the Gaussian at infinity, the indicator of \(y_0\) at the places dividing \(\mathfrak m\), and the vacuum at the other finite places.  For a vector \(F\) on \(\Hom(E,E^{12})\) write \(\Jet(F)\) for its jet series at scale \(\pi_M\cjet\); for \(F\) with the Gaussian at infinity, as the datum of \(y_0\), this is the Taylor series at \(0\), in the twelve copies, of \(z \mapsto \Th_F(\pi_M\cjet z)\), as in Definition~\ref{d2-def:jet-series}.  Write \(\Jet_{y_0}\) for this series for the datum of \(y_0\).  Two facts are read from the proof.  First, if \(3e_{i,\beta} \le c_\beta\) for all \(i \le 12\) and all \(\beta\), the product over the twelve copies of the series \(\prod_\beta y(\mu_\beta z)^{e_{i,\beta}}\,\tth(\mu_\beta z)^{c_\beta}\) is \(\eta\) applied to \(\Jet(F)\) for a vector \(F\) on \(\Hom(E, E^{12})\) whose values at finite arguments are integral at \(\eta\) and whose theta function is entire and fixed by the signed translations along \(\pi_M\mathfrak a_0^{12}\), with \(\mathfrak a_0\) as in Definition~\ref{def:twelve}.  Second, for such \(F\), finite Fourier inversion writes \(M^{12f}F\) as a finite sum \(\sum_j c_j\cdot(\text{translate of the vacuum})\) such that \(\lvert\eta(c_jv_*)\rvert \le 1\) for every value \(v_*\) of the vacuum at a finite argument, and each translate is a sum over \(y_0\) of roots of unity times the datum of \(y_0\).  \emph{Status}: both are read from the proof; for background on algebraic theta functions compare Bannai and Kobayashi, Duke Math.\ J.\ 153 (2010).  Here the integral vacuum of \(\Hom(E,E^{12})\) is \(\cvac^{12}\) times the vacuum (Section~\ref{ob-sec:vacuum}).  Also read from the proof: every conjugate under \(\Aut(\C)\) of a value of the integral vacuum of \(\Hom(E, E^{12})\) at a finite argument has absolute value at most \(\exp(v_0q_1^2)^{12}\).

\begin{problem}\label{pr:one-column}
Let \(y_1, \dots, y_N \in \C_p\), \(c_1, \dots, c_N\) with \(\lvert c_i\rvert \le 1\), and \(\lvert\sum_i c_iy_i\rvert \ge \epsilon_0 > 0\).  Then \(\lvert y_i\rvert \ge \epsilon_0\) for some \(i\).  Fix a modulus as above and a value \(c_*\) of the integral vacuum of \(\Hom(E,E^{12})\) at a finite argument (with \(\val(0) = +\infty\)).  If \(\lv_{12}(\lambda) \le \rho\) and \(c_\beta \ge 3\lvert\lambda\rvert\) for all \(\beta\), there are \(\hb\) with \(\sum_i h_{\tau,i} = \lambda(\tau)\) and \(y_0\) with \(\val\eta\bigl(\cvac^{12}[z^{\hb}]\Jet_{y_0}\bigr) \le \rho + \val\eta(M^{12f}c_*)\).
\end{problem}

Fix a nonzero value \(c_*\) of the integral vacuum of \(\Hom(E, E^{12})\) at a finite argument, before any modulus or prime is chosen, and put \(\catt := M^{12f}\,c_*\prod_\tau \Pi_\tau^2\).  If \(\lambda(\tau) \le \kappa(\tau) + x^{42}\) for every \(\tau\), then \(\lvert\lambda\rvert \le f(b + x^{42}) \le 2fb\), so \(c_\beta \ge 6fb \ge 3\lvert\lambda\rvert\).

\begin{problem}\label{pr:c-att}
If \(\lambda(\tau) \le \kappa(\tau) + x^{42}\) for every \(\tau\), then for every \(\sigma \in \Aut(\C)\)
\[
  \log\lvert\sigma(\catt)\rvert \le 12f\log M + 2fu\log(2b + u) + 12v_0q_1^2 \le \frac{fmb}{k} .
\]
\end{problem}

Problems~\ref{pr:monomial-witness}, \ref{pr:one-column} and~\ref{pr:c-att}, with the bound on twelve copies (to which Problems~\ref{pr:partners}, \ref{pr:six-unit}, \ref{pr:near-unique} and~\ref{pr:loss-product} contribute), the tight frame and the two facts on theta functions, prove the following theorem.

\begin{theorem}\label{thm:attained}
Let \(\lambda\colon\Phi\to\N\) satisfy \(\kappa(\tau) \le \lambda(\tau) \le \kappa(\tau) + x^{42}\) for every \(\tau\).  There is \(\catt \in \C\), \(\catt \ne 0\), with \(\lvert\sigma(\catt)\rvert \le \exp(fmb/k)\) for every \(\sigma \in \Aut(\C)\), such that: for every modulus \(\mathfrak m\) of \(E\) with ideal \((\pi_M)\), every prime \(p\) and every ring homomorphism \(\eta\colon\C\to\C_p\), there are a multi-index \(\hb = (h_{\tau,i})\) on the twelve copies with \(\sum_{i} h_{\tau,i} = \lambda(\tau)\) for every \(\tau\), and a class \(y_0 \in (\Oc_E/\mathfrak m)^{12}\), with
\[
  \val\,\eta\Bigl(\cvac^{12}\,\catt^{-1}\prod_{\tau\in\Phi}\lambda(\tau)!\cdot[z^{\hb}]\,\Jet_{y_0}\Bigr) \le \stot .
\]
\end{theorem}

For fixed \(\lambda\) the same multiplier \(\catt\) works for all \(\mathfrak m\), \(p\) and \(\eta\); the witnesses \(\hb\), \(y_0\) may depend on all three.  Chapter~\ref{l5-chap} applies the theorem with \(\lambda\) the row \(\lambda_P\) of a critical pair main for \(j\) (Definitions~\ref{l3-def:row} and~\ref{l3-def:main}), which satisfies the hypothesis by Problem~\ref{l3-pr:row-range}.

Taken as given: the normal form, the isotropic scalings, the tight frame, the six-copy span and coupling bounds, the twelve-copy bound, the vacuum bound and the two facts on theta functions.  The next sheet needs Theorem~\ref{thm:attained} and \eqref{eq:stot}.

\paragraph{In the Lean formalization.}
In the Lean proof \cite{GL} (readable edition, \texttt{BoundAttained.lean}), Theorem~\ref{thm:attained} is \texttt{Data.exists\_rpow\_neg\_localPrimeSumTotal\_le\_norm\_coeff}, stated with norms, \(\Jet_{y_0}\) being \texttt{twelveTaylor} of \texttt{indicatorDatum} and \(\stot\) \texttt{Data.localPrimeSumTotal}.  Definitions~\ref{def:lv} and~\ref{def:near}: \texttt{CMLattice.LeastValLE}, \texttt{Data.IsNearPair}.  Normal form: \texttt{Data.exists\_hasNormalForm\_resolutionForms}.  Isotropic scalings: \texttt{CMLattice.expSubst\_sub\_rescale\_mem}, \texttt{CMLattice.prod\_rescale\_theta\_mem}.  Tight frame: \texttt{Data.exists\_isTightFrame}.  Coupling bound: \texttt{CMLattice.leastValLE\_partnerDegree\_of\_isCouplingFamily}.  Twelve copies: \texttt{Data.leastValLE\_twelve}.  The first fact on theta functions is \texttt{Data.exists\_integral\_isThetaFunction\_twelveTaylor\_eq}; the second is the Fourier expansion (\texttt{Space.exists\_fourier\_coefficients}) in the proof of \texttt{Data.exists\_norm\_le\_norm\_indicatorDatum}.  The coefficient assertion in Problem~\ref{pr:one-column} is \texttt{Data.exists\_le\_norm\_coeff\_indicatorDatum}.  Further steps: \texttt{Data.mul\_sum\_weight\_sub\_localPrimeSumTotal}, \texttt{Parameters.cmCurve}, \texttt{Parameters.periodFactor}, \texttt{Parameters.resolutionForms}, \texttt{CMLattice.leastValLE\_of\_unit\_coefficients}, \texttt{Data.leastValLE\_six\_zero\_of\_inert}, \texttt{Data.isCouplingFamily\_of\_isNearPair}, \texttt{CMLattice.LeastValLE.exists\_monomial}, \texttt{Parameters.norm\_ringEquiv\_intVacuum\_le} and \texttt{Data.norm\_map\_mul\_intVacuum\_le}.

\section*{Hints}

\subsection*{\S\ref{sec:lemniscate}}
\begin{description}
\item[\ref{pr:lemniscate-least}] Problem~\ref{ba-pr:monomial}, by machine.
\item[\ref{pr:coupling-exp}] Expand the exponential; for the integral series use Problem~\ref{ba-pr:product} and \(\mu(6) \ge E_3(6)\).
\item[\ref{pr:one-partner}] By machine: the series of \(F = \sigma^2\wp\) through degree \(52\) from \(\wp = \sll^{-2}\) and \((\log\sigma)'' = -\wp\) (Section~\ref{st-sec:moments}), then \(F(s+t)F(\zeta(s-t))\).
\end{description}

\subsection*{\S\ref{sec:partners}}
\begin{description}
\item[\ref{pr:partners}] Compare term by term.
\item[\ref{pr:coupled-pair}] Euler's criterion.
\end{description}

\subsection*{\S\ref{sec:six}}
\begin{description}
\item[\ref{pr:isotropic}] Hensel's lemma at \(13\); for general \(p\), Lagrange's theorem and the compactness of \(\Z_p^4\).
\item[\ref{pr:unit-inverse}] Invert the one-variable product in \(\oo[[y]]\), then substitute the forms.
\item[\ref{pr:six-unit}] Theorem~\ref{ba-thm:katz} in total degree \(\lvert\nu\rvert\) along \(\mu_{\beta_0}\), then specialize the six copies along an isotropic vector with first entry \(1\).
\end{description}

\subsection*{\S\ref{sec:near}}
\begin{description}
\item[\ref{pr:near-unique}] Start with (b): \(\Phi_b\) lies in one class modulo \(q_0\); compute \(d^e\) modulo \(q_0\), using \(d \equiv q_0/2 - 1\) and \(4 \mid q_0\).  For (a) at an inert prime, Problem~\ref{ba-pr:near-pairs} for \(\Phi\), \([-p]\), \(\Phi_b\), \(L = l_0(s)\), the sparsity in Definition~\ref{da-def:data} and Problem~\ref{dl-pr:qone-split}.
\end{description}

\subsection*{\S\ref{sec:twelve}}
\begin{description}
\item[\ref{pr:monomial-witness}] As in Problem~\ref{ba-pr:monomial}: truncate, and use the ultrametric inequality.
\item[\ref{pr:one-column}] Separate the case \(c_* = 0\).  Otherwise apply the first part to the Fourier expansion multiplied by the vacuum value \(c_*/\cvac^{12}\), then to the expansion of a translate over the \(y_0\); Problem~\ref{pr:monomial-witness} for the exponents.
\item[\ref{pr:c-att}] Bound each factor of \(\Pi_\tau\) by \(2b + u\), and \(\lvert\sigma(c_*)\rvert\) by the bound stated before Problem~\ref{pr:one-column}; for the last inequality use \(M \le b^{22}\), \(k \le r\) and the size conditions of Definition~\ref{da-def:data}.
\end{description}

\begin{thebibliography}{GL}
\bibitem[GL]{GL} L.~Sernau, \emph{GRH for Dirichlet \(L\)-functions, in Lean~4}, a shorter rewriting of the Lean proof, text by Claude, \url{https://github.com/lsernau/grh-lean}, commit \texttt{72fead1}, readable edition \texttt{grh/}.
\end{thebibliography}

\end{document}
